% Section 5: Discussion
% Target: 2 pages (~1200-1400 words)
% Focus: Societal impact, generalizability, limitations, ethical considerations

\section{Discussion}
\label{sec:discussion}

\subsection{Implications for Cultural Preservation}
\label{subsec:cultural_impact}

The 10.5\% improvement in cultural authenticity demonstrates that AI systems can support rather than undermine cultural diversity. Applications include: (1) \textbf{Language Education} for diaspora communities via culturally-grounded explanations in learning platforms; (2) \textbf{Digital Archiving} providing systematic documentation while translation enables accessibility to researchers and policy-makers; (3) \textbf{Cultural Connection} through future interactive chatbot interfaces enabling conversational exploration of proverbs.

\subsection{Generalizability and Scalability}
\label{subsec:generalizability}

While evaluated on Kikuyu, the methodology generalizes to: (1) Related Bantu languages (Kikamba, Kimeru) sharing linguistic structures; (2) Other African languages (Yoruba, Amharic) with rich proverbial traditions; (3) Global indigenous languages (Māori, Quechua, Navajo) facing comparable challenges. Partnership with Africa Proverbs Working Group provides pathways for multi-language expansion. The \ograg{} architecture extends beyond proverbs to folktales (narrative pattern ontologies), ceremonial discourse (pragmatic context), and multimodal heritage (linking textual, visual, auditory artifacts).

\subsection{Limitations and Future Validation}
\label{subsec:limitations}

\textbf{Evaluation Methodology}: The most significant limitation is reliance on automated evaluation rather than formal human judgment. Future work should recruit 15-20 native speakers representing dialectical, generational, and geographic diversity for community-based validation with structured rubrics (target inter-rater reliability Krippendorff's $\alpha \geq 0.67$).

\textbf{Ontology Coverage}: Current coverage is 60-70\% of documented concepts. Error analysis revealed 68\% of failures trace to missing concepts. Targeted expansion prioritizing 15-20 high-frequency concepts could increase coverage to 85-90\%.

\textbf{Cultural Diversity}: The ontology represents ``standard'' Kikuyu from elder perspectives, but cultural knowledge varies generationally, dialectically, and contextually. Representing this internal diversity while maintaining ontological tractability remains an open challenge.

\subsection{Ethical Considerations}
\label{subsec:ethics}

Cultural knowledge is intellectual property requiring consent and benefit-sharing. Sustainable deployment should transfer cultural authority to Kikuyu organizations through joint ownership, community veto power, revenue sharing, and training programs. User interfaces should emphasize epistemic humility. Risks include cultural ossification (formalizing dynamic traditions), misappropriation (exploitation without community benefit), and oversimplification. Mitigation strategies: versioning, licensing restrictions, and transparency documentation.

\subsection{Theoretical Contributions}
\label{subsec:theory}

\\textbf{The Structure Hypothesis}: Structured knowledge representations outperform unstructured ones when cultural concepts exhibit relational organization. The 5.3\\% gap between \\ograg{} and \\trad{} (0.627 vs. 0.584) empirically validates this, suggesting domains with rich conceptual interrelationships benefit from graph-based vs. vector-based retrieval.
